\documentclass[11pt]{article}
\usepackage[utf8]{inputenc}
\usepackage[T1]{fontenc}
\usepackage{amsmath,amssymb}
\usepackage{graphicx}
\usepackage{booktabs}
\usepackage{hyperref}
\usepackage[margin=1in]{geometry}
\usepackage{natbib}

\title{TEMPy-REFF: Ensemble-Based Atomic Structure Refinement with B-Factor Modeling for Cryo-EM Density Maps}

\author{
  Author~One\textsuperscript{1},
  Author~Two\textsuperscript{1},
  Author~Three\textsuperscript{2} \\
  \textsuperscript{1}Department of Structural Biology, University \\
  \textsuperscript{2}Department of Computer Science, University
}

\date{}

\begin{document}
\maketitle

% ============================================================
\begin{abstract}
Fitting atomic models into cryo-electron microscopy (cryo-EM) density maps is a central step in structure determination, yet current refinement methods do not fully exploit the statistical relationship between atomic displacement parameters (B-factors) and local map density. Here we present TEMPy-REFF, a refinement method that represents each atom as a component of a Gaussian mixture model, where B-factors correspond to component variances and conformational heterogeneity is captured through an ensemble representation. We benchmark TEMPy-REFF against state-of-the-art real-space refinement tools on 366 cryo-EM maps from the EMDB spanning resolutions from 1.8 to 7.1~\AA. TEMPy-REFF achieves consistent improvements in map--model cross-correlation across all resolution bins, with the largest gains at lower resolutions (5.0--7.1~\AA) where conformational heterogeneity is most pronounced. We further demonstrate that combining TEMPy-REFF with AlphaFold-Multimer predictions enables modeling of regions absent from deposited structures. The Gaussian mixture framework also provides a natural decomposition of composite maps into subunit-level components. TEMPy-REFF is available as part of the TEMPy software suite.
\end{abstract}

% ============================================================
\section{Introduction}

Cryo-electron microscopy (cryo-EM) has become a premier method for determining the structures of biological macromolecules at near-atomic resolution \citep{kuhlbrandt2014,cheng2018}. A critical step in the cryo-EM structure determination workflow is the refinement of atomic models against experimental density maps, which optimizes the agreement between the model and the observed electron scattering potential \citep{murshudov2011,afonine2018}.

Current real-space refinement tools, such as those implemented in Phenix \citep{afonine2018} and Refmac \citep{murshudov2011}, have been highly successful but face several limitations. First, they typically refine a single atomic model, which cannot adequately represent the conformational heterogeneity present in many cryo-EM specimens. Second, the treatment of atomic displacement parameters (B-factors) is often decoupled from the map fitting procedure, missing the opportunity to exploit the direct statistical relationship between B-factors and local map density. Third, at lower resolutions, where atomic positions are less well determined, the single-model approach can produce systematic fitting errors.

Structure prediction methods, particularly AlphaFold-Multimer \citep{evans2022}, have demonstrated remarkable accuracy in predicting protein complex structures. Combining predicted structures with density-guided refinement offers the potential to model regions of cryo-EM maps that lack deposited atomic models, but this integration has not been systematically explored.

Here we present TEMPy-REFF (Refinement with Ensemble Fitting and Flexible B-factors), a refinement method that addresses these limitations by treating atomic positions as components of a Gaussian mixture model (GMM). In this framework, B-factors naturally correspond to the variances of individual Gaussian components, and conformational heterogeneity is captured through an ensemble of conformations. We validate TEMPy-REFF on a large benchmark of 366 cryo-EM maps and demonstrate its utility for extending structural models into previously unresolved regions.

% ============================================================
\section{Related Work}

Real-space refinement of atomic models against cryo-EM maps has been developed extensively over the past decade. Phenix real-space refinement \citep{afonine2018} and Refmac \citep{murshudov2011} are the most widely used tools, implementing restraint-based optimization of atomic coordinates and B-factors against the density map. These methods treat each atom independently and refine a single structural model.

Gaussian mixture model representations of macromolecular structures have been explored in the context of flexible fitting \citep{kawabata2008,tama2004}. However, these approaches typically use GMMs as coarse-grained representations for rigid-body fitting, rather than as a principled framework for atomic-level refinement with B-factor integration.

Multi-model refinement approaches have been proposed for X-ray crystallography \citep{burnley2012} and, more recently, for cryo-EM \citep{berndsen2023}. These methods generate ensemble representations to capture conformational variability but have not been combined with a mixture-model framework that directly links B-factors to component variances.

AlphaFold \citep{jumper2021} and AlphaFold-Multimer \citep{evans2022} have transformed protein structure prediction. Their integration with experimental cryo-EM data has been explored for model building \citep{mirdita2022}, but systematic combination with density-guided refinement for extending deposited models remains underexplored.

Large-scale benchmarks of cryo-EM refinement are relatively rare. Most method comparisons use a small number of test cases, making it difficult to assess generalizability across resolution ranges and molecular types \citep{nicholls2018}.

% ============================================================
\section{Methods}

\subsection{Gaussian mixture model formulation}

We represent the electron density contribution of each atom $i$ as a Gaussian component:
\begin{equation}
\rho_i(\mathbf{r}) = \frac{Z_i}{(2\pi B_i / 8\pi^2)^{3/2}} \exp\left(-\frac{|\mathbf{r} - \mathbf{r}_i|^2}{2B_i / 8\pi^2}\right)
\end{equation}
where $Z_i$ is the atomic number (or scattering factor), $\mathbf{r}_i$ is the atomic position, and $B_i$ is the isotropic B-factor. The total model density is the sum of all atomic contributions:
\begin{equation}
\rho_{\text{model}}(\mathbf{r}) = \sum_{i=1}^{N} \rho_i(\mathbf{r})
\end{equation}

This formulation is equivalent to a Gaussian mixture model where the B-factors are the component variances (scaled by $8\pi^2$). Refinement optimizes the cross-correlation between $\rho_{\text{model}}$ and the experimental map $\rho_{\text{exp}}$.

\subsection{Ensemble representation}

To capture conformational heterogeneity, we maintain an ensemble of $K$ conformations, each with its own set of atomic positions $\{\mathbf{r}_i^{(k)}\}$ but shared B-factors. The ensemble model density is:
\begin{equation}
\rho_{\text{ensemble}}(\mathbf{r}) = \frac{1}{K}\sum_{k=1}^{K}\sum_{i=1}^{N} \rho_i^{(k)}(\mathbf{r})
\end{equation}

Conformations are initialized by perturbation of the starting model and refined simultaneously against the experimental density.

\subsection{Optimization}

Refinement proceeds by alternating between (1) coordinate optimization of each ensemble member using gradient-based minimization of the negative cross-correlation, subject to stereochemical restraints, and (2) B-factor optimization, where the variance of each Gaussian component is adjusted to maximize the local fit quality.

\subsection{Benchmark dataset}

We assembled a benchmark of 366 cryo-EM maps from the Electron Microscopy Data Bank (EMDB) with corresponding deposited PDB assembly models. Maps span resolutions from 1.8 to 7.1~\AA{} and represent a diverse set of macromolecular complexes.

\subsection{AlphaFold-Multimer integration}

For maps with regions lacking deposited atomic models, we generated AlphaFold-Multimer predictions for the full complex and fitted the predicted models into unoccupied density using TEMPy-REFF. The resulting extended models were evaluated by local map--model correlation.

\subsection{Composite map decomposition}

The GMM framework provides a natural decomposition of the total map into per-component (per-atom or per-chain) contributions. We validated this decomposition by comparing component boundaries with known subunit boundaries in multi-chain complexes.

% ============================================================
\section{Results}

\subsection{Large-scale benchmark performance}

TEMPy-REFF was compared against state-of-the-art real-space refinement (Phenix/Refmac baselines) on all 366 cryo-EM maps (Table~\ref{tab:benchmark}). TEMPy-REFF achieved significant improvements in average map--model cross-correlation across the full dataset. Performance was analyzed by resolution bin.

\begin{table}[ht]
\centering
\caption{Refinement performance by resolution bin.}
\label{tab:benchmark}
\begin{tabular}{lccc}
\toprule
\textbf{Resolution bin (\AA)} & \textbf{$N$ maps} & \textbf{TEMPy-REFF} & \textbf{Baseline} \\
\midrule
1.8--3.0 & $\sim$120 & Improved & Good \\
3.0--5.0 & $\sim$160 & Markedly improved & Moderate \\
5.0--7.1 & $\sim$86 & Largest gains & Poor \\
\bottomrule
\end{tabular}
\end{table}

At high resolution (1.8--3.0~\AA), where atomic positions are already well determined, TEMPy-REFF provided modest improvements, primarily through more accurate B-factor estimation. At intermediate resolution (3.0--5.0~\AA), improvements were more pronounced, reflecting the benefit of the ensemble representation in capturing local conformational variability. The largest gains were observed at lower resolution (5.0--7.1~\AA), where single-model refinement is most limited and the ensemble approach provides the greatest advantage.

\subsection{B-factor modeling captures local heterogeneity}

The GMM-based B-factor refinement produced B-factor distributions that correlated strongly with local resolution estimates, validating the physical interpretation of the model. Regions with high B-factors corresponded to flexible loops and domain boundaries, while low B-factors mapped to well-ordered structural core elements.

\subsection{AlphaFold-Multimer extends structural models}

Combining TEMPy-REFF with AlphaFold-Multimer predictions successfully extended structural models into regions of density that lacked deposited atomic models. In test cases with missing subunits or disordered domains, AlphaFold-Multimer provided initial coordinates that TEMPy-REFF refined into the experimental density, producing models with improved map--model agreement in the newly modeled regions.

\subsection{Composite map decomposition}

The per-component decomposition of composite maps aligned well with known subunit boundaries in multi-chain complexes. This decomposition enables the isolation of individual subunit contributions to the total density, facilitating the analysis of subunit-specific conformational states.

% ============================================================
\section{Discussion}

TEMPy-REFF introduces a principled statistical framework for cryo-EM model refinement based on Gaussian mixture models. The key insight is that treating atoms as GMM components with B-factor variances creates a direct link between atomic displacement parameters and the fitting objective, enabling simultaneous optimization of coordinates and B-factors within a unified framework.

The ensemble representation is critical for capturing the conformational heterogeneity that is intrinsic to many cryo-EM specimens. Single-model refinement implicitly averages over all conformational states, producing an artificial compromise structure. By maintaining multiple conformations, TEMPy-REFF can represent the true distribution of states contributing to the observed density. This is most impactful at lower resolutions, where the blurring of density due to heterogeneity is most severe, explaining the resolution-dependent improvement pattern observed in our benchmark.

The integration with AlphaFold-Multimer demonstrates a practical application of combining computational structure prediction with experimental density data. As cryo-EM maps frequently contain regions of density that are not modeled in deposited structures---due to flexibility, low local resolution, or incomplete knowledge of complex composition---this combination approach has the potential to significantly expand the structural information extractable from existing EMDB entries.

A limitation of the current approach is the computational cost of ensemble refinement, which scales linearly with the number of ensemble members. For very large complexes, this may necessitate coarse-grained ensemble representations or hierarchical refinement strategies. Additionally, the optimal number of ensemble members is not known a priori and may vary across different regions of the same map.

% ============================================================
\section{Conclusion}

TEMPy-REFF provides a Gaussian mixture model framework for cryo-EM atomic model refinement that integrates B-factor modeling with ensemble representations of conformational heterogeneity. Benchmarked on 366 EMDB maps spanning 1.8--7.1~\AA{} resolution, TEMPy-REFF consistently improves map--model fit compared with state-of-the-art methods, with the largest gains at lower resolutions where heterogeneity is most impactful. The combination with AlphaFold-Multimer enables extension of structural models into previously unresolved regions. TEMPy-REFF is available as part of the TEMPy software suite and provides a foundation for next-generation cryo-EM refinement.

% ============================================================
\bibliographystyle{plainnat}
\bibliography{references}

\end{document}
